%% Determines the type of document (including standard settings for layout).
\documentclass[letterpaper]{article}



%% Package to control the size of a page: the standard margins used by LaTeX are
%% very wide!
\usepackage[margin=1in]{geometry}


%% Packages add functionality to the document. The AMS packages are standard
%% packages to support various mathematical notations. AMS stands for American
%% Mathematical Society, the organization that maintains these packages.
\usepackage{amsmath,amsthm,amssymb}

%% Theorems-like environments using functionality provided by amsthm.
\theoremstyle{plain}
\newtheorem{theorem}{Theorem}[section]
    %% [section] at the end specifies that theorems should be numbered
    %% per-section: Section x starts with theorem-like x.1 and so on, ...
\newtheorem{proposition}[theorem]{Proposition}
    %% [theorem] in the middle specifies: use the same counter as the theorem
    %% environment: here we number all theorem-like environments consecutively.
\newtheorem{corollary}[theorem]{Corollary}
\newtheorem{lemma}[theorem]{Lemma}
\theoremstyle{definition}
\newtheorem{definition}[theorem]{Definition}
\theoremstyle{remark}
\newtheorem{example}[theorem]{Example}
\newtheorem{remark}[theorem]{Remark}


%% Support for nicely formatted tables.
\usepackage{booktabs}


%% Support for colors & colors in tables.
\usepackage[table]{xcolor}

% Seven colors safe for use color blindness.
% Colors taken from doi:10.1038/nmeth.1618.
\definecolor{cbOrange}{RGB}{230,159,0}
\definecolor{cbSkyBlue}{RGB}{86,180,233}
\definecolor{cbBluishGreen}{RGB}{0,158,115}
\definecolor{cbYellow}{RGB}{240,228,66}
\definecolor{cbBlue}{RGB}{0,114,178}
\definecolor{cbVermillion}{RGB}{213,94,0}
\definecolor{cbReddischPurple}{RGB}{204,121,167}


%% Notation used in this document.
\newcommand{\n}{\mathbf{n}} %% Num. Replicas.
\newcommand{\f}{\mathbf{f}} %% Num. Faulty Replicas.

%% Misc. Math notation.
\newcommand{\BigO}{\mathcal{O}}
\newcommand{\abs}[1]{\lvert #1 \rvert}
\newcommand{\AName}[1]{\textsc{#1}}
\newcommand{\Var}[1]{\texttt{#1}}


%% Algorithms.
\usepackage{algorithm}
\usepackage[noend]{algorithmic}
\newcommand{\GETS}{:=}


%% Formatting SI-units.
\usepackage{siunitx}
\sisetup{per-mode=symbol}


%% TikZ: for creating figures.
\usepackage{tikz}

%% Configuration for figures: Nicer arrows.
\usetikzlibrary{arrows.meta}
\tikzset{>=Stealth}


%% pgfplots: drawing plots using TikZ.
\usepackage{pgfplots}
%% Configuration for plots: Use color-blind friendly colors.
\pgfplotscreateplotcyclelist{cbSafeList}{
    very thick,solid,cbOrange,every mark/.append style={solid},mark=*\\
    very thick,solid,cbSkyBlue,every mark/.append style={solid},mark=*\\
    very thick,solid,cbBluishGreen,every mark/.append style={solid},mark=*\\
    very thick,solid,cbYellow,every mark/.append style={solid},mark=*\\
    very thick,solid,cbBlue,every mark/.append style={solid},mark=*\\
    very thick,solid,cbVermillion,every mark/.append style={solid},mark=*\\
    very thick,solid,cbReddischPurple,every mark/.append style={solid},mark=*\\
    very thick,solid,black,every mark/.append style={solid},mark=*\\
}
\pgfplotsset{
    legend style={font=\small},
    compat=1.16,
    width=260pt,
    height=140pt,
    legend cell align=left,
    xlabel near ticks,
    ylabel near ticks,
    every axis/.append style={
        cycle list name=cbSafeList,
        ymin=0,
        enlargelimits=0.05,
        mark size=1pt,
        ylabel style={align=center},
        xlabel style={align=center},
        title style={align=center}
    }
}

%% PgfplotsTable: loading data files to use with pgfplots.
\usepackage{pgfplotstable}


%% Support for hyperlinks and urls. The setting ``colorlinks'' sets how links
%% are shown in the document (with a color, without underline). We put hyperref
%% last---it has a tendency to break other packages when loaded before them.
\usepackage[colorlinks]{hyperref}
\usepackage{graphicx}
\usepackage{pgffor}
\usepackage{caption}
\captionsetup[algorithm]{labelformat=empty, labelsep=none}
\usepackage{tabularx}
\usepackage{enumitem}
\usepackage{fancyhdr}
\usepackage{comment}



\begin{document}

%% First the meta-data of the doucment: title,

\begin{titlepage}
    \centering
    \vspace*{3cm}
    
    {\Huge COMPSCI 2DB3 Assignment 7}\\[2cm]
    
    {\large Luca Mawyin}\\[0.5cm]
    
    {\large Dr. Jelle Hellings}\\[0.5cm]

    {\large \today}\\[2cm]

    {\Large McMaster University}

    \vfill
    \begin{flushleft}

        {\large \textbf{Student Number: }400531739}\\[0.5cm]
        {\large \textbf{MacID: }mawyinl}
    \end{flushleft}

\end{titlepage}

\renewcommand{\thesection}{Part \Roman{section}}
\renewcommand{\thesubsection}{Problem \arabic{section}}
\renewcommand{\thesubsubsection}{P\arabic{section}.\arabic{subsubsection}}

\section{}
\subsection{}
\begin{align*}
    &\mathbf{R}(a,b,c,d,e,f)\\
    &\mathbb{D}=\{a\to b,c,~c,d\to e,~d\to f,~f\to a,c,f\}
\end{align*}

\begin{algorithm}
    \textbf{Algorithm} CLOSURE($\mathbb{D},~A$):
    \begin{algorithmic}[1]
        \STATE{\textit{result}~\GETS~$A$}
        \WHILE{there exists a $X\to Y\in\mathbb{D}$ with $X\subseteq \textit{result}$ and $Y\nsubseteq \textit{result}$}
            \STATE{\textit{result}~\GETS~$\textit{result}\cup Y$}
        \ENDWHILE{}
        \STATE{\textbf{end while}}
        \RETURN{\textit{result}}
    \end{algorithmic}
\end{algorithm}

\subsubsection{}
\begin{enumerate}
    \item $X=\{a,d\}$
    \item $\textit{result}=\{a,d\}$
    \item Use $a\to b,c$, as $a\in\textit{result}\land{Y\nsubseteq\textit{result}}$
    \item $\textit{result}=\{a,b,c,d\}$
    \item Use $c,d\to e$, as $c,d\in\textit{result}\land{Y\nsubseteq\textit{result}}$
    \item $\textit{result}=\{a,b,c,d,e\}$
    \item Use $d\to f$, as $d\in\textit{result}\land{Y\nsubseteq\textit{result}}$
    \item $\textit{result}=\{a,b,c,d,e,f\}$
    \item We have, among other things, $a,d\to f$
\end{enumerate}
\subsubsection{}
\begin{enumerate}
    \item $a\to b,c$
    \item $d\to a,b,c,d,e,f$
    \item $f\to a,b,c,f$
    \item $a,f\to a,b,c,f$
\end{enumerate}
Only the function dependencies that meaningully rely on each other were counted. This means if key $\{b,c\}$ does not result in any dependencies, then it is not included. Furthermore, any pair that does not meaningfully contribute to another key is not included. For example, in pair $\{a,e\}$, as $e$ does not help in the dependency, the dependency is the same as $a$, and subsequently the pair is not included.

As $d$ is a superkey, any set containing $d$ is also a superkey. This includes, but is not limited to sets such as $\{a,d\}$, $\{b,d\}$, $\{c,d\}$, etc. As $d$ is the only single attribute that can be used to identify all other attributes, this means that $d$ and only $d$ is a candidate key. Furthermore, as $d$ is required for attribute $e$, and no attribute $x$ nor set $X$ exists such that $x\to d$ or $X\to d$, any set that does not include $d$ will not be a superkey. 
\newpage
\subsubsection{}
\begin{algorithm}
    \textbf{Algorithm} MINIMALCOVER($\mathbb{D}$):
    \begin{algorithmic}[1]
        \STATE{\textit{result}~\GETS~$\{X\to y~|~(X\to Y\in\mathbb{D})\land (y\in Y)\}$}
        \WHILE{$\mathbb{D}$ is not a minimal cover}
            \IF{there exists a $X\to y\in \textit{result}$ that is redundant w.r.t. \textit{result}}
                \STATE{\textit{result}~\GETS~$\textit{result}\backslash \{X\to y\}$}
            \ELSIF{\{there exists a $X\to y\in \textit{result}$ with $a\in X$ superfluous w.r.t. \textit{result}\}}
                \STATE{\textit{result}~\GETS~$(\textit{result}\backslash\{X\to y\})\cup\{X\backslash a \to y\}$}
            \ENDIF{}
            \STATE{\textbf{end if}}
        \ENDWHILE{}
        \STATE{\textbf{end while}}
        \RETURN{\textit{result}}
    \end{algorithmic}
\end{algorithm}
\begin{enumerate}
    \item $\textit{result}=\{a\to b;~a\to c;~c,d\to e;~d\to f;~f\to a;~f\to c;~f\to f\}$
    \item $f \to f$ is redundant, therefore remove
    \item $\textit{result}=\{a\to b;~a\to c;~c,d\to e;~d\to f;~f\to a;~f\to c\}$
    \item $f \to c$ is redundant as $f\to a$ and $a\to c$, therefore remove
    \item $\textit{result}=\{a\to b;~a\to c;~c,d\to e;~d\to f;~f\to a;\}$
    \item $c,d\to e$ is not redundant, therefore test if $c\to e$ or $d\to e$
    \item As neither exists, we are left with $\textit{result}=\{a\to b;~a\to c;~c,d\to e;~d\to f;~f\to a\}$
    \item As no other dependency is redundant, and each untested $X$ is a single attribute, each LHS is not redundant
    \item We return $\textit{result}=\mathbb{D}'=\{a\to b;~a\to c;~c,d\to e;~d\to f;~f\to a\}$
\end{enumerate}
\subsubsection{}
The relational schema \textbf{R} is not in 3NF. This is because to be in 3NF, every FD $A\to B\in\mathbb{D}^*$, $A$ must be a candidate key of \textbf{R}. However, the only candidate key in \textbf{R} is $d$, meaning there are other FDs where $A$ is not a candidate key, subsequently breaking the rule.\\[1em]
Performing the for-loop at the start of the algorithm we get:
\begin{align*}
    \{a\}&: \textbf{S}_1(\underline{a},b,c);\\
    \{d\}&: \textbf{S}_2(\underline{d},f);\\
    \{f\}&: \textbf{S}_3(\underline{f},a);\\
    \{c,d\}&: \textbf{S}_4(\underline{c},\underline{d},e);\\
\end{align*}
The condition of the while-loop does not hold for any pair of relations in $\textbf{S}_1,\ldots,\textbf{S}_4$. Furthermore, the condition of the if-statement does not hold. Therefore, we return: 
\begin{align*}
    \{a\}&: \textbf{S}_1(\underline{a},b,c);\\
    \{d\}&: \textbf{S}_2(\underline{d},f);\\
    \{f\}&: \textbf{S}_3(\underline{f},a);\\
    \{c,d\}&: \textbf{S}_4(\underline{c},\underline{d},e);\\
\end{align*}
\newpage
\noindent For each \textbf{S} we have the dependencies as follows:
\begin{enumerate}[label=$\textbf{S}_\text{\arabic*}$]
    \item $a\to b,c$
    \item $d\to f$
    \item $f\to a$
    \item $c,d\to e$
\end{enumerate}
This decomposition is a lossless-join, as we can recover the original dependencies easily, and there were no additional dependencies added. Furthermore, the decomposition is dependency preserving, as every dependency appears in at least one relation.
\subsubsection{}
The relational schema \textbf{R} is not in BCNF. This is because to be in BCNF, every FD $A\to B\in\mathbb{D}^*$, $A$ must be a primary key of \textbf{R}. However, there exist some FDs such as $a\to b,c$, and as $a$ is not a primary key, this breaks the rule.\\[1em]
First we initialize:
\begin{align*}
    \textit{result}=\{\textbf{R}\}\\
\end{align*}
Next, we perform the while-loop. The functional dependency $a\to b,c\in\mathbb{D}^*$ is a BCNF violation for \textbf{R}. Therefore we replace \textbf{R} by two relations $\textbf{S}_1$ and $\textbf{S}_2$:
\begin{align*}
    \textit{result}=\{&\textbf{S}_1(\underline{a},b,c);\\
    &\textbf{S}_2(a,\underline{d},e,f)\}
\end{align*}
In $\textbf{S}_2$, the functional dependency $f\to a$ violates BCNF. Therefore we replace $\textbf{S}_2$ by two relations $\textbf{S}_{2,1}$ and $\textbf{S}_{2,2}$:
\begin{align*}
    \textit{result}=\{&\textbf{S}_1(\underline{a},b,c);\\
    &\textbf{S}_{2,1}=(\underline{f},a);\\
    &\textbf{S}_{2,2}(\underline{d},e,f)\}
\end{align*}
The decomposition is a lossless join, as each step of decomposition uses a dependency where the determinant is a key in one of the resulting relations. As such, each split is a lossless join. Furthermore, the decomposition is not dependency preserving, as we lose some dependencies that were implicit, such as $f\to b,c$ as $f\to a$ and $a\to b,c$. However, this decomposition cannot enforce the implicit dependency, and therefore is not dependency preserving.
\subsubsection{}
The relational schema \textbf{R} is not in 4NF . This is because for any FD $A\to B\in\mathbb{D}^*$ that holds on the schema \textbf{R}, it must satisfy $A\twoheadrightarrow B$. However, as $b$ is not a superkey, this violates 4NF.\\[1em]
First, we initialize:
\begin{align*}
    \textit{result}=\{\textbf{R}\}\\
\end{align*}
Next, we perform the while-loop. the multivalued dependency $b\twoheadrightarrow c\in\mathbb{D}^*$ is a 4NF violation. We replace \textbf{R} by two relations $\textbf{S}_1$ and $\textbf{S}_2$:
\begin{align*}
    \textit{result}=\{&\textbf{S}_1(\underline{b},c)\\
    &\textbf{S}_2(a,b,\underline{d},e,f)\}
\end{align*}
In $\textbf{S}_2$, $a\to b$ violates 4NF. We replace $\textbf{S}_2$ by two relations $\textbf{S}_{1,1}$ and $\textbf{S}_{2,2}$:
\begin{align*}
    \textit{result}=\{&\textbf{S}_1(\underline{b},c)\\
    &\textbf{S}_{2,1}(\underline{a},b)\\
    &\textbf{S}_{2,2}(a,\underline{d},e,f)\}
\end{align*}
In $\textbf{S}_{2,2}$, $f\to a$ violates 4NF. We replace $\textbf{S}_{2,2}$ by two relations $\textbf{S}_{2,2,1}$ and $\textbf{S}_{2,2,2}$:
\begin{align*}
    \textit{result}=\{&\textbf{S}_1(\underline{b},c)\\
    &\textbf{S}_{2,1}(\underline{a},b)\\
    &\textbf{S}_{2,2,1}(a,\underline{f})\\
    &\textbf{S}_{2,2,2}(\underline{d},e,f)\}
\end{align*}
\section{}
\subsection{}
\subsubsection{}
\begin{align*}
    \mathbb{D}'=\{\textit{isbn}&\to \textit{imprint};\\
    \textit{isbn}&\to \textit{btitle};\\
    \textit{isbn}&\to \textit{year};\\
    \textit{isbn}&\to \textit{edition};\\
    \textit{isbn}&\to \textit{numrecipes};\\
    \textit{imprint}&\to \textit{publisher};\\
    \textit{imprint,btitle,year}&\to \textit{isbn};\\
    \textit{imprint,btitle,edition}&\to \textit{isbn};\\
    \textit{rid}&\to \textit{isbn};\\
    \textit{rid}&\to \textit{rtitle};\\
    \textit{rid}&\to \textit{rpage}\\
    \textit{isbn,rtitle}&\to \textit{rid};\\
    \textit{aid}&\to \textit{aname}\}
\end{align*}
The relational schema \textbf{BookInfo} is not in 3NF. This is because to be in 3NF, every FD $A\to B\in\mathbb{D}^*$, $A$ must be a candidate key of \textbf{BookInfo}. However, this is not the case.\\[1em]
Performing the for-loop at the start of the algorithm we get:
\begin{align*}
    \textit{isbn}&: \textbf{S}_1 (\textit{\underline{isbn},imprint,btitle,year,edition,numrecipes});\\
    \textit{imprint}&: \textbf{S}_2 (\textit{\underline{imprint},publisher});\\
    \textit{imprint,btitle,year}&: (\textbf{S}_3 \textit{\underline{imprint},\underline{btitle},\underline{year},isbn});\\
    \textit{imprint,btitle,edition}&: \textbf{S}_4 (\textit{\underline{imprint},\underline{btitle},\underline{{edition}},isbn});\\
    \textit{rid}&: \textbf{S}_5 (\textit{\underline{rid},isbn,rtitle,rpage});\\
    \textit{isbn,rtitle}&: \textbf{S}_6 (\underline{isbn},\underline{rtitle},rid);\\
    \textit{aid}&: \textbf{S}_7 (\underline{aid},aname)
\end{align*}
After the while loop we will remove $\textbf{S}_6$, $\textbf{S}_3$, and $\textbf{S}_4$:
\begin{align*}
    \textit{isbn}&: \textbf{S}_1 (\textit{\underline{isbn},imprint,btitle,year,edition,numrecipes});\\
    \textit{imprint}&: \textbf{S}_2 (\textit{\underline{imprint},publisher});\\
    \textit{rid}&: \textbf{S}_5 (\textit{\underline{rid},isbn,rtitle,rpage});\\
    \textit{aid}&: \textbf{S}_7 (\underline{aid},aname)
\end{align*}
Next we will add the relation
\begin{align*}
    \text{candidate key}&:\textbf{S}_8(\underline{isbn},\underline{rid},\underline{aid})
\end{align*}
This decomposition is a lossless-join as there exists a relation which contains a key of the original schema. As we have candidate key $\textbf{S}_8$, this holds true. Furthermore, the decomposition is dependency preserving, as we maintain the original dependencies of the relational schema.
\subsubsection{}
The relational schema \textbf{BookInfo} is not in BCNF. This is because to be in BCNF, every FD $A\to B\in\mathbb{D}^*$, $A$ must be a primary key of \textbf{BookInfo}. However, this does not hold, as we have some keys that are not primary keys that still exist within out schema.\\[1em]
First we initialize:
\begin{align*}
    \textit{result}&=\{\textbf{BookInfo}\}
\end{align*}
Next, we perform the while loop. The FD $\textit{isbn}\to\textit{publisher,imprint,btitle,year,edition,numrecipes}\in\mathbb{D}^*$ is a BCNF violation for \textbf{BookInfo}. Hence, we replace \textbf{BookInfo} by two relations $\textbf{S}_1$ and $\textbf{S}_2$:
\begin{align*}
    \textit{result}=\{&\textbf{S}_1(\textit{\underline{isbn},publisher,imprint,btitle,year,edition,numrecipes}),\\
    &\textbf{S}_2(\textit{isbn,\underline{rid},rtitle,rpage,\underline{aid},aname})\}
\end{align*}
In $\textbf{S}_2$ we have BCNF violation $\textit{rid}\to\textit{isbn,rtitle,rpage}\in\mathbb{D}^*$. Therefore we replace $\textbf{S}_2$ by relations $\textbf{S}_{2,1}$ and $\textbf{S}_{2,2}$:
\begin{align*}
    \textit{result}=\{&\textbf{S}_1(\textit{\underline{isbn},publisher,imprint,btitle,year,edition,numrecipes}),\\
    &\textbf{S}_{2,1}(\textit{\underline{rid},isbn,rtitle,rpage});\textbf{S}_{2,2}(\textit{\underline{isbn},\underline{aid},aname})\}
\end{align*}
In $\textbf{S}_{2,2}$ we have violation $\textit{aid}\to\textit{aname}\in\mathbb{D}^*$. We will replace $\textbf{S}_{2,2}$ by $\textbf{S}_{2,2,1}$ and $\textbf{S}_{2,2,2}$:
\begin{align*}
    \textit{result}=\{&\textbf{S}_1(\textit{\underline{isbn},publisher,imprint,btitle,year,edition,numrecipes}),\\
    &\textbf{S}_{2,1}(\textit{\underline{rid},isbn,rtitle,rpage});\\
    &\textbf{S}_{2,2,1}(\textit{\underline{aid},aname})\\
    &\textbf{S}_{2,2,2}(\textit{\underline{isbn},aid})\}
\end{align*}
This decomposition is lossless-join, as every step performed is lossless. This is because for some arbitrary relations $\textbf{S}_1$ and $\textbf{S}_2$, we will have $\textbf{S}_1\cap\textbf{S}_2=X$, where $X$ is our original schema. Furthermore, the decomposition also dependency-preserving, as we can enforce every FD in $\mathbb{D}^*$ using only the decomposed relations
\end{document}
